% Part: first-order-logic
% Chapter: completeness
% Section: introduction

\documentclass[../../../include/open-logic-section]{subfiles}

\begin{document}

\iftag{FOL}
      {\olfileid{fol}{com}{int}}
      {\olfileid{pl}{com}{int}}

\olsection{سريزه}

د بشپړتيا قضيه د منطق په اړه له تر ټولو بنسټيزو پايلو څخه يوه ده۔
دا په دوو بڼو وړاندې کېږي، او موږ به د دواړو همارزښت ثابت کړو۔ په
لومړۍ بڼه کښې دا زموږ د معنٰی پوهنې د پايلې او د !!{derivation}
نظام ترمنځ د اړيکې په اړه يو بنسټيز حقيقت وايي: که
!!a{sentence}~$!A$ د ځينو !!{sentence}s~$\Gamma$ معنايي پايله
وي، نو داسې !!a{derivation} هم شته چې
$\Gamma \Proves !A$ ثابتوي۔ په دې توګه، د !!{derivation} نظام
هغومره پياوړے دے څومره چې د هغو خبرو له ثابتولو پرته کېدے شي چې
په رښتيا نۀ لازمېږي۔

په دويمه بڼه کښې دا د مدل د شتون د پايلې په توګه ويل کېدے شي: د
!!{sentence}s هر سازګار سټ د صدق وړ دے۔ سازګاري يو ثبوت-تيوريکي
مفهوم دے: دا وايي چې زموږ د !!{derivation} نظام د يو ټاکلي ډول
!!{derivation}s نۀ شي جوړولے۔ خو دا څوک تضمينوي چې يوازې ځکه چې
له~$\Gamma$ څخه د يو ټاکلي ډول هېڅ !!{derivation}s نشته، خامخا
داسې \iftag{FOL}{!!a{structure}~$\Struct{M}$}{!!{valuation}{~$\pAssign{v}$}
شته چې $\iftag{FOL}{\Sat{M}}{\pSat{v}}{\Gamma}$}؟ د بشپړتيا د
قضيې تر لومړي ثبوت وړاندې---بلکې د هغو !!{derivation} نظامونو تر
رامنځته کېدو هم وړاندې چې اوس ئې لرو---ستر جرمن رياضي‌پوه ډېوېډ
هيلبرټ دا نظر درلود چې د رياضي نظريو سازګاري د هغو څيزونو شتون
تضمينوي چې نظريې ئې په اړه دي۔ هغه ګوتلوب فرېګه ته په يو ليک کښې
داسې وويل:
\begin{quote}
  که په خپل‌سري ډول ورکړل شوي بديهي اصول له خپلو ټولو پايلو سره يو
  بل نۀ نقضوي، نو رښتوني دي او هغه څيزونه چې بديهي اصول ئې تعريفوي
  شتون لري۔ زما دپاره همدا د رښتياوالي او شتون معيار دے۔
\end{quote}
فرېګه سخت مخالفت وکړ۔ د بشپړتيا د قضيې دويمه بڼه ښيي چې هيلبرټ
لږ تر لږه په دې معنٰی سم و چې که بديهي اصول سازګار وي، نو داسې
\emph{کوم}
\iftag{FOL}{!!{structure}}{!!{valuation}} شته چې ټول رښتوني کوي۔

دا يوازيني لاملونه نۀ دي چې د بشپړتيا قضيه---يا لا کره، د هغې
ثبوت---مهمه ده۔ دا يو شمېر مهمې پايلې لري چې ځينې به ئې جلا وڅېړو۔
د بېلګې په توګه، ځکه هر !!{derivation} چې
$\Gamma \Proves !A$ ښيي متناهي دے، نو له~$\Gamma$ څخه يوازې
متناهي ډېر !!{sentence}s کارولے شي۔ له دې امله د بشپړتيا له قضيې
څخه لازمېږي چې که~$!A$ د~$\Gamma$ معنايي پايله وي، نو لا له وړاندې
د~$\Gamma$ د کوم متناهي فرعي سټ معنايي پايله ده۔ دې ته
\emph{د متناهي شاهد خاصيت} ويل کېږي۔ په همارز ډول، که د~$\Gamma$
هر متناهي فرعي سټ سازګار وي، نو په خپله~$\Gamma$ هم سازګار دے۔

که څه هم د متناهي شاهد قضيه د بشپړتيا له قضيې څخه د
!!{derivation}s د منځنۍ لارې په تېرېدو لازمېږي، د بشپړتيا د قضيې
\emph{خپله ثبوت} هم نېغ په نېغه د هغې د ثابتولو دپاره کارول کېدے
شي۔ ځکه ثبوت د !!{sentence}s داسې سټ اخلي چې يو ټاکلے خاصيت---
سازګاري---لري، او له دې سټه داسې !!a{structure} جوړوي چې ټاکلي
خاصيتونه لري (په دې حالت کښې، سټ پوره کوي)۔ کابو همدا جوړونه د
متناهي شاهد خاصيت د نېغ ثبوت دپاره کارېږي، که د سازګارو سټونو پر
ځاے د !!{sentence}s له «په متناهي ډول د صدق وړ» سټونو څخه پيل
وکړو۔
\iftag{FOL}{دا جوړونه نورې پايلې هم ورکوي؛ د بېلګې په توګه، د
  !!{sentence}s هر د صدق وړ سټ يو متناهي يا !!{denumerable}s مدل
  لري۔ (دې پايلې ته د لوېنهايم--سکولم قضيه ويل کېږي۔) په عمومي ډول،
  د !!{sentence}s له سټونو څخه د !!{structure}s جوړول په منطق کښې
  ډېر ځله، او کله کله حتٰی په فلسفه کښې هم، کارول کېږي۔}{}

\end{document}
